The High-Luminosity (HL) upgrade \cite{hl-lhc-tdr} of the Large Hadron Collider (LHC) \cite{lhc-tdr} will substantially increase its instantaneous luminosity, leading to a greater number of collisions per bunch crossing. The HL-LHC is expected to reach an instantaneous luminosity of up to $7.5 \times 10^{34} \, \text{cm}^{-2}\text{s}^{-1}$, resulting in an average of up to 200 simultaneous proton-proton interactions per bunch crossing, commonly referred to as \textit{pileup}. Pileup denotes additional interactions occurring in the same bunch crossing as the primary interaction of interest. In comparison, in Run 3 the average pileup reached up to about 60 per bunch crossing. This elevated pileup environment poses a major challenge for the ATLAS experiment \cite{atlas}, as it complicates the reconstruction of particle tracks and the identification of individual collisions.

To meet the challenges posed by high pileup conditions, the ATLAS collaboration has launched a major upgrade program, including the complete replacement of the inner detector with a new all-silicon tracking system, the Inner Tracker (ITk) \cite{itk-pixel-tdr}\cite{itk-strip-tdr}. In parallel, the reconstruction software must also be enhanced to cope with the increased event complexity. For track reconstruction, the ATLAS collaboration has decided to adopt the Acts (a common tracking software) toolkit \cite{acts-2022} --- a community-driven project that provides experiment-independent tracking algorithms written in modern C++ \cite{atlas-hl-lhc-computing-cdr}. By leveraging Acts, ATLAS benefits from state-of-the-art developments in tracking algorithms and software architecture, while ensuring flexibility and adaptability for future requirements.

This thesis strongly contributes to the ongoing effort of improving Acts and its performance for the HL-LHC conditions, and to integrate it into the ATLAS tracking software.

The document is divided into three main parts: introduction, developments, and performance. The first part provides a general introduction to the field of high-energy physics, the reconstruction of tracks and primary vertices, and Acts. The second part focuses on the developments made to Acts related to several areas: stepper and navigator, track finding, tracking in dense material, and tracking with time information. The third part presents the tracking performance of the OpenDataDetector with Acts and the performance of the Acts-based track reconstruction for ATLAS ITk. Summaries, conclusions, and outlooks are provided on a chapter or part level.
